\chapter{Tutorial 4: aiming at one band}\label{ch:tut-review}

Suppose that a weak observed band could be a dark band of the calculation,
but the computed intensity does not match. This tutorial shows how to
recompute only what concerns that band, more accurately, without
repeating the whole calculation.
\begin{enumerate}
\item In \file{resonancecandidates}, find the quantum-number changes that
  involve the dark state. Lines marked \texttt{near} are close to a
  threshold. Set their last column to \texttt{on}.
\item In \file{bandlist}, leave active only the band of interest (and the
  bright band it borrows from), setting the others to \texttt{off}.
\item Rerun phase 1 with
\begin{lstlisting}
resonanceoverrides = "resonancecandidates"
bandlistfile = "bandlist"
\end{lstlisting}
  then phases 2a and 2b, only for the starting states that matter.
\end{enumerate}
% TODO: worked example on pyrene with numbers.
